% Options for packages loaded elsewhere
\PassOptionsToPackage{unicode}{hyperref}
\PassOptionsToPackage{hyphens}{url}
%
\documentclass[
]{article}
\usepackage{amsmath,amssymb}
\usepackage{iftex}
\ifPDFTeX
  \usepackage[T1]{fontenc}
  \usepackage[utf8]{inputenc}
  \usepackage{textcomp} % provide euro and other symbols
\else % if luatex or xetex
  \usepackage{unicode-math} % this also loads fontspec
  \defaultfontfeatures{Scale=MatchLowercase}
  \defaultfontfeatures[\rmfamily]{Ligatures=TeX,Scale=1}
\fi
\usepackage{lmodern}
\ifPDFTeX\else
  % xetex/luatex font selection
\fi
% Use upquote if available, for straight quotes in verbatim environments
\IfFileExists{upquote.sty}{\usepackage{upquote}}{}
\IfFileExists{microtype.sty}{% use microtype if available
  \usepackage[]{microtype}
  \UseMicrotypeSet[protrusion]{basicmath} % disable protrusion for tt fonts
}{}
\makeatletter
\@ifundefined{KOMAClassName}{% if non-KOMA class
  \IfFileExists{parskip.sty}{%
    \usepackage{parskip}
  }{% else
    \setlength{\parindent}{0pt}
    \setlength{\parskip}{6pt plus 2pt minus 1pt}}
}{% if KOMA class
  \KOMAoptions{parskip=half}}
\makeatother
\usepackage{xcolor}
\usepackage[margin=1in]{geometry}
\usepackage{graphicx}
\makeatletter
\def\maxwidth{\ifdim\Gin@nat@width>\linewidth\linewidth\else\Gin@nat@width\fi}
\def\maxheight{\ifdim\Gin@nat@height>\textheight\textheight\else\Gin@nat@height\fi}
\makeatother
% Scale images if necessary, so that they will not overflow the page
% margins by default, and it is still possible to overwrite the defaults
% using explicit options in \includegraphics[width, height, ...]{}
\setkeys{Gin}{width=\maxwidth,height=\maxheight,keepaspectratio}
% Set default figure placement to htbp
\makeatletter
\def\fps@figure{htbp}
\makeatother
\setlength{\emergencystretch}{3em} % prevent overfull lines
\providecommand{\tightlist}{%
  \setlength{\itemsep}{0pt}\setlength{\parskip}{0pt}}
\setcounter{secnumdepth}{-\maxdimen} % remove section numbering
\usepackage{custom}
\ifLuaTeX
  \usepackage{selnolig}  % disable illegal ligatures
\fi
\IfFileExists{bookmark.sty}{\usepackage{bookmark}}{\usepackage{hyperref}}
\IfFileExists{xurl.sty}{\usepackage{xurl}}{} % add URL line breaks if available
\urlstyle{same}
\hypersetup{
  pdftitle={STAT 581 - Exam 1: Due Dec 2, 2021},
  pdfauthor={Alex Towell (atowell@siue.edu)},
  hidelinks,
  pdfcreator={LaTeX via pandoc}}

\title{STAT 581 - Exam 1: Due Dec 2, 2021}
\author{Alex Towell
(\href{mailto:atowell@siue.edu}{\nolinkurl{atowell@siue.edu}})}
\date{}

\begin{document}
\maketitle

\hypertarget{problem-1}{%
\section{Problem 1}\label{problem-1}}

\begin{quote}
An experiment is conducted to study the effect of fitness level on ego
strength. Random samples of college faculty members are selected from
each fitness level, and an ego score is observed for each member in the
sample. Higher values indicate greater ego. The data is provided as an
attachment.
\end{quote}

\vspace{50pt}

\hypertarget{a}{%
\subsection{(a)}\label{a}}

\begin{quote}
State the hypotheses of interest. Provide an interpretation, stated in
the context of the problem.
\end{quote}

\vspace{50pt}

\hypertarget{b}{%
\subsection{(b)}\label{b}}

\begin{quote}
Compute the \(t_0\) statistic and the \(p\)-value. Provide an
interpretation, statedin the context of the problem.
\end{quote}

\vspace{50pt}

\hypertarget{c}{%
\subsection{(c)}\label{c}}

\begin{quote}
Create a Boxplot as a graphical display of the data. Is it true that all
high fitness faculty members have greater egos than low fitness faculty
members?
\end{quote}

\vspace{50pt}

\hypertarget{d}{%
\subsection{(d)}\label{d}}

\begin{quote}
Compute a 95\% confidence interval for \(\delta = \mu_1 - \mu_2\).
Provide an interpretation, stated in the context of the problem.
\end{quote}

\vspace{50pt}

\hypertarget{e}{%
\subsection{(e)}\label{e}}

\begin{quote}
Explain how a confidence interval provides a complementary result to a
hypothesis test.
\end{quote}

\vspace{50pt}

\hypertarget{f}{%
\subsection{(f)}\label{f}}

\begin{quote}
Explain how a confidence interval can be used in testing
\(H_0 : \mu_1 = \mu_2\).
\end{quote}

\vspace{50pt}

\hypertarget{problem-2}{%
\section{Problem 2}\label{problem-2}}

\begin{quote}
A completely randomized design is used to investigate the effect of drug
dosage on the activity level of lab rats. Each dose level is applied to
\(n=4\) rats, and an activity score is observed for each rat in the
sample. Higher values indicate greater activity. The data is provided as
an attachment.
\end{quote}

\vspace{50pt}

\hypertarget{a-1}{%
\subsection{(a)}\label{a-1}}

\begin{quote}
State the statistical hypotheses of interest. Briefly explain how the
form of the alternative hypothesis requires a need for further
investigation.
\end{quote}

\vspace{50pt}

\hypertarget{b-1}{%
\subsection{(b)}\label{b-1}}

\begin{quote}
Compute the \(F_0\) statistic and the p-value. Provide an
interpretation, stated in the context of the problem. Create a Boxplot
as a graphical display of the data.
\end{quote}

\vspace{50pt}

\hypertarget{c-1}{%
\subsection{(c)}\label{c-1}}

\begin{quote}
Compute and display 95\% confidence intervals for all pairwise
comparisons. Explain how computing multiple intervals impacts the
probability of committing an error.
\end{quote}

\vspace{50pt}

\hypertarget{d-1}{%
\subsection{(d)}\label{d-1}}

\begin{quote}
Perform pairwise comparisons using the Fisher LSD method, and the Tukey
method. Provide grouping information for each method. Comment on the
seemingly contradictory nature of a pairwise comparisons analysis.
\end{quote}

\vspace{50pt}

\hypertarget{e-1}{%
\subsection{(e)}\label{e-1}}

\begin{quote}
Describe the defining characteristics for each of the above pairwise
comparison methods.
\end{quote}

\vspace{50pt}

\hypertarget{f-1}{%
\subsection{(f)}\label{f-1}}

\begin{quote}
Compute the margin of error and comparison-wise error rate for the Tukey
method in this problem.
\end{quote}

\vspace{50pt}

\hypertarget{g}{%
\subsection{(g)}\label{g}}

\begin{quote}
Compute the margin of error and family-wise error rate for the Fisher
LSD method in this problem.
\end{quote}

\vspace{50pt}

\hypertarget{problem-3}{%
\section{Problem 3}\label{problem-3}}

\begin{quote}
A factorial experiment is used to investigate the effect of pressure,
temperature, and time on the yield from a chemical reaction. Two levels
(\texttt{low}, \texttt{high}) of each factor are set and \(n=2\) runs of
a \(2^3\) design are completed. The data is provided as an attachment.
\end{quote}

\vspace{50pt}

\hypertarget{a-2}{%
\subsection{(a)}\label{a-2}}

\begin{quote}
Perform tests for all main effects and for all interaction effects.
State the \(F\)-statistic and \(p\)-value for each test of an effect
deemed to be important. Fit a reduced model with the main effects and
the statistically significant interaction effect.
\end{quote}

\vspace{50pt}

\hypertarget{b-2}{%
\subsection{(b)}\label{b-2}}

\begin{quote}
Provide a general definition of an interaction effect. Explain how an
interaction plot is used in studying an interaction effect.
\end{quote}

\vspace{50pt}

\hypertarget{c-2}{%
\subsection{(c)}\label{c-2}}

\begin{quote}
Create a plot for the interaction effect deemed important. Provide an
interpretation, stated in the context of the problem.
\end{quote}

\vspace{50pt}

\hypertarget{d-2}{%
\subsection{(d)}\label{d-2}}

\begin{quote}
Create a Boxplot showing the main effect for the remaining factor.
Provide an interpretation, stated in the context of the problem.
\end{quote}

\vspace{50pt}

\hypertarget{e-2}{%
\subsection{(e)}\label{e-2}}

\begin{quote}
Create a plot of the fitted values for the reduced model. Which setting
of the factors should be used if the goal is to maximize yield?
\end{quote}

\vspace{50pt}

\hypertarget{f-2}{%
\subsection{(f)}\label{f-2}}

\begin{quote}
Explain how the analysis is providing a simplification to the observed
data.
\end{quote}

\vspace{50pt}

\hypertarget{problem-4}{%
\section{Problem 4}\label{problem-4}}

\begin{quote}
An experiment to compare a new drug to a standard is in the planning
stages. The response variable of interest is the clotting time (in
minutes) of blood drawn from the subject. The experimenters want to
perform a two sample \(t\) test at level \(\alpha = .05\) with power
\(\pi = .90\) at \(\delta_A = 0.5\), for standard deviation
\(\sigma = .7\).
\end{quote}

\vspace{50pt}

\hypertarget{a-3}{%
\subsection{(a)}\label{a-3}}

\begin{quote}
Determine the sample size for each drug in order to achieve the stated
test specifications.
\end{quote}

\vspace{50pt}

\hypertarget{b-3}{%
\subsection{(b)}\label{b-3}}

\begin{quote}
Graph the power curve for the chosen sample size. Explain how the power
curve displays the desired properties of the test.
\end{quote}

\vspace{50pt}

\hypertarget{c-3}{%
\subsection{(c)}\label{c-3}}

\begin{quote}
Provide a general explanation of how A can be determined.
\end{quote}

\vspace{50pt}

\hypertarget{d-3}{%
\subsection{(d)}\label{d-3}}

\begin{quote}
Briefly discuss some other issues that may provide additional insight to
an experimental result beyond a finding of statistical significance.
\end{quote}

\vspace{50pt}

\hypertarget{e-3}{%
\subsection{(e)}\label{e-3}}

\begin{quote}
Briefly comment on the additional information provided by the p-value,
beyond a determination of statistical significance alone.
\end{quote}

\end{document}
